\section{Présentation du projet et de mon rôle}

Le projet IPF 2026, d'abord intitulé \textit{5 Secondes Chrono} puis décliné sous le nom Lynx Immo dans l'application, a été réalisé pour Le Carré Pro, organisme de formation spécialisé dans l'immobilier commercial. Le besoin initial consistait à proposer aux professionnels du secteur un format de micro-apprentissage juridique, accessible depuis une application mobile et fondé sur des quiz courts et réguliers. Au fil des échanges, ce besoin s'est élargi : le produit devait également permettre au client d'administrer les contenus, de suivre la progression des apprenants, de gérer plusieurs niveaux d'abonnement et de préparer des fonctions de coaching plus avancées.

Le projet s'est déroulé de novembre 2025 à juillet 2026 avec une équipe de cinq étudiants issus de deux promotions. J'y ai occupé une double fonction de chef de projet et de référent technique. Cette position m'a placé à l'interface entre le client, l'organisation de l'équipe et les décisions d'architecture. Je devais transformer les attentes en une trajectoire réalisable, suivre l'avancement, arbitrer le périmètre, accompagner les membres de l'équipe et intervenir directement lorsque des choix techniques structuraient la suite du projet.

Lynx Immo est ainsi moins présenté dans ce dossier comme une démonstration isolée de Flutter, React ou Supabase que comme une expérience de conduite de projet. Le principal enjeu n'a pas été de produire une fonctionnalité particulière, mais de maintenir une cohérence entre un besoin client ambitieux, des ressources limitées, plusieurs composants techniques et une équipe dont les disponibilités évoluaient au cours de l'année.

\section{Compétences principalement mobilisées}

\begin{table}[H]
\centering
\begin{tabular}{p{2.5cm} p{11.5cm}}
\toprule
\textbf{Compétences} & \textbf{Éléments de preuve principaux} \\
\midrule
C1.4 à C1.7 & Feuille de route, jalons, risques, exigences, Jira, indicateurs, suivi des écarts, arbitrages et communication avec le client. \\
C4.1 à C4.5 & Identification des compétences, définition des rôles, répartition et adaptation de la charge, accompagnement et montée en compétences. \\
C4.6 & Rétrospective collective et entretiens individuels de fin de projet, dans les limites d'un contexte académique sans processus RH. \\
C3.6 & Documentation développeur MkDocs et organisation technique destinées à permettre la reprise du produit. \\
\bottomrule
\end{tabular}
\caption{Compétences principalement mobilisées dans le projet Lynx Immo}
\label{tab:competences-lynx}
\end{table}

Le contexte académique impose une limite importante pour C4.2 et C4.6. L'équipe était déjà constituée et aucun service de ressources humaines n'intervenait dans le projet. Je n'ai donc pas conduit de recrutement ni établi de plans de carrière. En revanche, j'ai participé à la constitution opérationnelle de l'équipe : identification des compétences nécessaires, attribution des responsabilités, intégration progressive des membres, accompagnement et évaluation du fonctionnement collectif et individuel.

\section{Passer d'une idée de quiz à un produit structuré}

Lors des premiers échanges, le produit pouvait encore être résumé comme une application de quiz juridiques. Cette définition est rapidement devenue insuffisante. Pour proposer un service réellement exploitable, il fallait relier l'expérience apprenant à un espace d'administration, à une base de données commune, à des mécanismes de progression, à des contenus vidéo, à des statistiques et à une logique d'abonnement. Chaque ajout paraissait cohérent pris isolément, mais leur accumulation risquait de rendre le projet irréalisable dans le temps disponible.

La première décision de pilotage a donc été de distinguer le produit cible du produit effectivement livrable. Avec l'équipe, nous avons défini un parcours principal : un utilisateur devait pouvoir créer son compte, renseigner son profil, accéder à un quiz quotidien, répondre aux questions, consulter ses résultats et progresser. En parallèle, le client devait disposer d'un espace lui permettant de gérer les contenus indispensables à ce parcours. Les fonctions de messagerie, de coaching avancé ou d'analyse plus poussée restaient importantes, mais ne devaient pas empêcher la réalisation d'un socle cohérent.

Ce cadrage s'est matérialisé dans plusieurs documents : cahier des charges, Project Scope Statement, Plan de Développement Logiciel, exigences fonctionnelles et non fonctionnelles, registre des risques, maquettes, modèle de données, roadmap et planning. Au début du projet, ce travail documentaire pouvait sembler retarder le développement. Il est pourtant devenu essentiel lorsque la trajectoire réelle s'est éloignée des hypothèses initiales : nous disposions alors d'une référence permettant de distinguer ce qui était indispensable, négociable ou reportable.

\subsection{Construire une feuille de route adaptée au calendrier réel}

La feuille de route couvrait la période de novembre 2025 à juillet 2026. Elle organisait le projet en grandes phases : cadrage, préparation des contenus, validation du design, développement du MVP, stabilisation, préparation des environnements, sécurité et finalisation. Un diagramme de Gantt et des sprints permettaient ensuite de décliner cette vision en travaux opérationnels.

Je n'ai pas cherché à appliquer un rythme Scrum strict indépendamment du contexte. Les deux promotions ne disposaient pas des mêmes périodes d'alternance, des mêmes rendus ni des mêmes disponibilités. Les sprints ont donc conservé une logique itérative, mais leur durée et leur contenu ont été adaptés aux contraintes scolaires et aux rendez-vous client. Ce choix me paraissait plus honnête qu'une cadence théorique qui aurait masqué la capacité réelle de l'équipe.

\begin{figure}[H]
\centering
\fbox{\parbox{0.9\textwidth}{\centering
Emplacement prévu : synthèse de la roadmap IPF et de ses principaux jalons.\\
La version détaillée du Gantt sera placée en annexe.}}
\caption{Trajectoire générale du projet Lynx Immo}
\label{fig:lynx-roadmap}
\end{figure}

La feuille de route n'a pas seulement servi à annoncer des dates. À chaque jalon, elle permettait de comparer le niveau de maturité atteint au niveau attendu, puis de décider de la suite. Lorsque le MVP s'est révélé moins avancé que prévu, conserver toutes les fonctionnalités aurait conduit à disperser davantage l'équipe. Le planning a alors rempli sa fonction principale : rendre l'écart visible et fournir un support à l'arbitrage. Cette démarche constitue ma preuve principale de C1.4.

\section{Organiser une équipe aux compétences et disponibilités différentes}

L'équipe réunissait cinq étudiants aux profils différents. J'assurais la direction du projet et la référence technique ; Johan Nedelec occupait le rôle de Product Owner ; Anne-Charlotte Arnold participait notamment au suivi opérationnel et à la documentation ; Clémence Yonkeu Wandji et William Plantec contribuaient au développement. Cette répartition initiale donnait un cadre, mais elle ne pouvait pas suffire à organiser plusieurs mois de travail.

Avant de distribuer les tâches, nous avons identifié les domaines nécessaires au projet : UX/UI avec Figma, développement mobile Flutter, administration React/Vite, Supabase et PostgreSQL pour le socle de données, sécurité, authentification, paiement, déploiement, documentation et pilotage. Cette cartographie a permis de vérifier que chaque domaine disposait d'un responsable ou d'un contributeur et de repérer les sujets nécessitant une montée en compétences. Elle matérialise C4.1 : l'équipe n'a pas été organisée uniquement autour des personnes disponibles, mais à partir des besoins de la solution.

Dans le cadre de C4.2, je distingue toutefois constitution juridique et constitution opérationnelle de l'équipe. Les membres n'ont pas été recrutés par mes soins. Mon travail a consisté à transformer un groupe déjà formé en une organisation capable de produire : préciser les rôles, déterminer les circuits de validation, associer les compétences aux lots de travail et revoir cette organisation lorsque les hypothèses de disponibilité n'étaient plus valables.

\subsection{Faire évoluer la répartition plutôt que défendre le plan initial}

Le projet a été découpé en grands domaines, puis en tâches opérationnelles suivies dans Jira. Au total, 563 tickets ont été créés pendant l'année. Ce nombre s'explique par l'étendue du produit : authentification, quiz, progression, badges, statistiques, administration, paiements, sécurité, données, environnements, documentation et préparation de la livraison. Jira permettait de relier la vision globale au travail quotidien, tandis que la matrice RACI précisait qui réalisait, validait, conseillait ou devait être informé.

La stratégie initiale prévoyait une forte implication des M2 dans le cadrage et les responsabilités transverses, puis une montée en puissance des M1 sur la réalisation. Le décalage d'un autre projet des M1 a perturbé cette organisation. Leur disponibilité est restée faible au moment où Lynx Immo devait accélérer son développement. Certains membres déjà engagés ont alors compensé une partie du retard, ce qui a permis au projet d'avancer mais a concentré la charge et la connaissance.

Cette situation m'a obligé à revoir les priorités et à redistribuer certains travaux. Elle a également montré la limite d'une répartition fondée uniquement sur les compétences : une personne peut être la plus pertinente techniquement pour une tâche sans disposer du temps nécessaire pour la prendre en charge. Coordonner l'équipe a donc consisté à arbitrer entre expertise, disponibilité, dépendances et progression des membres. C'est dans cette adaptation continue, plus que dans le RACI initial, que se situe ma principale démonstration de C4.3.

\section{Mesurer l'avancement et accepter ce que les indicateurs révélaient}

Pour éviter de piloter uniquement au ressenti, le projet associait plusieurs niveaux de suivi. Les exigences permettaient de mesurer la couverture du besoin ; Jira rendait visible l'état du travail ; les jalons comparaient la maturité obtenue à la trajectoire prévue ; le registre des risques signalait les menaces actives. Des KPI avaient également été définis pour le pilotage, la valeur d'usage et la performance technique.

Cette séparation était importante. Le nombre de fonctionnalités terminées ne suffisait pas à conclure qu'un produit apportait de la valeur, tout comme une fonctionnalité appréciée ne prouvait pas que le calendrier était maîtrisé. Les indicateurs de pilotage répondaient à la question de l'avancement ; les indicateurs d'usage auraient nécessité davantage d'utilisateurs réels et une période d'observation ; les indicateurs techniques préparaient le suivi des erreurs, de la disponibilité et des temps de réponse.

Le rapport final recensait 158 exigences, dont 67 fonctionnelles et 91 non fonctionnelles. Parmi elles, 104 étaient réalisées, 22 en cours et 32 non finalisées. Toutes les exigences classées urgentes avaient été couvertes, ce qui confirmait que l'arbitrage avait protégé le parcours principal sans permettre d'achever l'ensemble du produit cible.

\begin{figure}[H]
\centering
\fbox{\parbox{0.9\textwidth}{\centering
Emplacement prévu : état des exigences et vue synthétique du suivi Jira.}}
\caption{Suivi du projet par exigences et tâches Jira}
\label{fig:lynx-suivi}
\end{figure}

Le système de suivi n'a cependant pas été utilisé avec la même régularité pendant toute l'année. Autour de mars et avril, la charge et les retards ont réduit la mise à jour de certains documents. Le pilotage s'est alors recentré sur les jalons, les exigences prioritaires et les échanges directs. Cette limite est importante pour C1.6 : définir de nombreux KPI ne garantit pas leur utilité. Un indicateur n'a de valeur que s'il est alimenté, compris et relié à une décision. Sur un prochain projet, je conserverais moins d'indicateurs manuels et automatiserais davantage ceux qui peuvent être calculés depuis Jira, GitHub ou les outils de supervision.

\section{Superviser les écarts et arbitrer avec le client}

Le principal écart est apparu lorsque la disponibilité réelle de l'équipe a cessé de correspondre au planning. Les M2 ont continué à porter une part importante du cadrage, du pilotage et de la réalisation, tandis que plusieurs échéances académiques se concentraient sur la même période. Les retards ont alors eu un effet cumulatif : une tâche déplacée réduisait la marge du sprint suivant et augmentait la pression sur les personnes déjà fortement impliquées.

Le rôle du planning a évolué. Il ne s'agissait plus de respecter chaque date comme si les hypothèses initiales étaient toujours vraies, mais d'utiliser les jalons pour décider. Nous avons protégé le quiz quotidien, l'authentification, la progression, les badges, les statistiques principales et les fonctions d'administration indispensables. À l'inverse, la messagerie, les notifications, le coaching avancé et certains éléments de diffusion ont été limités ou reportés.

Les échanges réguliers avec le client ont facilité ces arbitrages, mais ils ont aussi fait émerger de nouvelles attentes. Je ne pouvais pas transformer chaque retour en obligation immédiate. Certaines demandes ont été intégrées lorsqu'elles corrigeaient un besoin mal compris ou amélioraient directement le MVP ; d'autres ont été simplifiées ou déplacées vers une version ultérieure. Le coaching illustre cette démarche : l'ambition d'un accompagnement fondé sur les données était pertinente, mais le projet ne disposait ni du volume d'usage ni du temps nécessaires à un système prédictif mature. Nous avons donc préparé le socle sans présenter cette partie comme achevée.

Cette confrontation entre le prévu, la capacité réelle et la valeur attendue constitue ma preuve principale de C1.7. Elle m'a appris qu'un planning utile n'est pas celui qui reste inchangé, mais celui qui permet de comprendre les conséquences d'un changement et d'en discuter avec les parties prenantes.

\section{Accompagner l'équipe sans concentrer toute la connaissance}

Plusieurs technologies étaient nouvelles pour certains membres. La montée en compétences a donc été pensée à partir des tâches : consultation de documentations officielles, tutoriels ciblés, prototypes, échanges internes et application progressive dans le produit. Le plan de formation couvrait notamment Supabase, Flutter, React, l'authentification, les paiements, la sécurité et le déploiement.

Dans les faits, cette montée en compétences n'a pas été homogène. Pendant les périodes tendues, la solution la plus rapide consistait souvent à confier une tâche complexe à la personne qui maîtrisait déjà le domaine. Ce choix répondait au besoin immédiat mais renforçait la dépendance envers quelques membres. La double fonction de chef de projet et de référent technique m'a moi-même placé dans cette situation : je pouvais relier rapidement une demande fonctionnelle à ses conséquences techniques, mais de nombreuses décisions et demandes d'aide remontaient alors vers la même personne.

L'accompagnement s'appuyait sur les réunions internes, Teams pour les échanges quotidiens, Jira pour le suivi, Confluence et les documents de cadrage pour les décisions, ainsi que GitHub pour la collaboration technique. Ces supports ont facilité la résolution de blocages et la transmission d'informations. Ils ne remplacent toutefois pas des temps de partage explicites. Avec le recul, j'organiserais davantage de binômage, de revues techniques courtes et de démonstrations internes. C4.4 et C4.5 sont donc démontrées par les dispositifs mis en place, mais aussi par l'analyse de leurs limites.

\section{Évaluer le fonctionnement de l'équipe}

L'évaluation a d'abord été collective. L'état des tâches, les exigences, les jalons, les risques et la maturité des livrables permettaient d'observer la capacité de l'équipe à produire. Une rétrospective de fin de projet a ensuite permis d'analyser ce qui avait bien fonctionné, la répartition du travail, l'utilité des outils et les améliorations possibles. Elle a confirmé la valeur du cadrage et des points d'avancement, tout en mettant en évidence une implication inégale et une charge importante sur certains membres.

Des entretiens individuels ont complété cette rétrospective. Ils portaient sur la place de chaque membre dans l'équipe, les tâches réalisées, les difficultés, le management et les améliorations envisageables. Leur objectif n'était pas de produire une notation descendante, mais de recueillir un retour réciproque et d'identifier les écarts entre la perception globale du projet et l'expérience individuelle.

Cette démarche fournit une preuve pour C4.6, tout en restant plus limitée qu'un dispositif d'entreprise. Il n'existait pas de référentiel individuel formel, de plan de carrière ni de service RH. Les entretiens ont en outre été organisés principalement en fin de projet. Ils démontrent néanmoins ma capacité à structurer un feedback, à analyser les contributions et le vécu des membres et à transformer ces retours en pistes d'amélioration pour une future équipe.

\section{Préparer une transmission réelle du produit}

La fin du projet ne devait pas se limiter à la soutenance académique. Le client peut décider de confier la suite à une entreprise ou à une nouvelle équipe ; il fallait donc éviter que le produit dépende de la mémoire de ses auteurs. Une documentation développeur MkDocs de plus d'une centaine de pages a été construite dans cette perspective. Elle couvre le démarrage, le périmètre, l'architecture, l'API, la base de données, le frontend, l'application mobile, le backend Supabase, les règles métier, la sécurité, les tests, le déploiement, la maintenance et les évolutions envisagées.

Le dépôt a également été préparé pour réduire les écarts d'environnement. Il regroupe le backend Supabase, l'administration React, l'application Flutter et les scripts de données. Un DevContainer fournit les principaux outils nécessaires ; le schéma de base est géré par migrations ; les procédures distinguent le local, la préproduction et la production et ajoutent des garde-fous avant les opérations destructrices. L'objectif est qu'une équipe de reprise puisse installer, comprendre et faire évoluer le projet sans dépendre de la configuration de mon poste.

\begin{figure}[H]
\centering
\fbox{\parbox{0.9\textwidth}{\centering
Emplacement prévu : page d'accueil et navigation de la documentation MkDocs.}}
\caption{Documentation développeur préparée pour la reprise de Lynx Immo}
\label{fig:lynx-mkdocs}
\end{figure}

Une réunion de transmission avec le client est prévue à la mi-août 2026. Au moment de la remise du présent dossier, cette réunion n'a pas encore eu lieu ; elle ne peut donc pas être présentée comme un résultat acquis. Elle constitue toutefois la dernière étape prévue de la démarche : présenter l'état réel du produit, les éléments fonctionnels, les limites connues, les procédures de lancement, la documentation et les conditions d'une reprise par un prestataire. Son compte rendu pourra servir de preuve complémentaire lors de la soutenance.

Cette préparation renforce C3.6, mais également C4.4 et C4.5. Dans le projet FCU, la documentation visait à rendre compréhensible un système historique. Dans Lynx Immo, elle vise à éviter qu'un produit récent devienne à son tour dépendant de quelques personnes. La documentation est ici un outil de continuité et de transmission, pas seulement un livrable de fin de projet.

\section{Bilan critique}

Lynx Immo est le projet qui m'a le plus confronté à la différence entre planifier et piloter. Produire un Gantt, définir des rôles ou préparer des KPI était relativement simple tant que les hypothèses restaient théoriques. La difficulté a commencé lorsque la disponibilité, le périmètre et le rythme de réalisation ont divergé du plan.

Je retiens d'abord qu'une feuille de route doit permettre de décider, et non seulement de constater un retard. Le planning a conservé sa valeur parce qu'il a rendu les écarts visibles et permis de protéger les exigences les plus importantes. Je retiens ensuite qu'un risque identifié n'est pas nécessairement maîtrisé. La disponibilité de l'équipe avait été repérée comme un risque, mais j'aurais dû réévaluer plus tôt la capacité réelle et réduire certaines ambitions avant que la charge ne se concentre.

Le troisième enseignement concerne la délégation. Confier un sujet à la personne la plus expérimentée accélère souvent la réalisation immédiate, mais peut fragiliser le collectif et la reprise. Sur un prochain projet, je consacrerais davantage de temps au binômage, aux revues et à la transmission progressive des domaines structurants. Enfin, la double responsabilité de chef de projet et de référent technique m'a permis d'arbitrer avec une compréhension précise des impacts, mais elle a également créé une dépendance et augmenté ma charge. Une séparation plus nette de ces rôles aurait facilité la délégation.

Le produit n'est pas présenté comme une application commerciale totalement industrialisée. Plusieurs fonctions restent à finaliser et une exploitation à grande échelle nécessiterait encore des tests utilisateurs, des validations de conformité et une phase de stabilisation supplémentaire. Le résultat est néanmoins une base fonctionnelle, versionnée et documentée, accompagnée d'un historique de décisions, d'un backlog et d'une démarche de transmission concrète.

Lynx Immo me permet ainsi de démontrer C1.4 à C1.7 ainsi que l'essentiel de C4.1 à C4.6. Pour C4.2 et C4.6, je conserve explicitement les limites du contexte académique : je n'ai pas conduit de recrutement ni de gestion de carrière. J'ai en revanche organisé une équipe autour de compétences identifiées, coordonné sa contribution, accompagné sa progression, évalué son fonctionnement et préparé la continuité du travail au-delà de l'équipe initiale.